\documentclass[12pt]{article}
\usepackage[T1]{fontenc}
\usepackage[utf8]{inputenc}
\usepackage[brazil]{babel}
\usepackage[margin=1in]{geometry}
\setlength{\parindent}{0pt}
\begin{document}
\section*{Tema 10}
Processo: PEDILEF 0504108- 62.2009.4.05.8200/ PB\\
Situacao: Julgado - TEMA 350/STF\\
Ramo: DIREITO PREVIDENCIÁRIO\\
Questao: Saber se há necessidade de renovação do requerimento administrativo para concessão de benefício assistencial.\\
Tese: O ajuizamento de ações visando o recebimento do benefício previsto no art. 203, V, da Constituição Federal não exige a renovação bienal do requerimento administrativo, afastada a indevida analogia ao art. 21, da Lei n. 8.742/93. RE 914797 - DECISÃO STF: determino a devolução dos autos ao Juízo de origem para que seja observada a decisão do Supremo no precedente RE 631.240 - TEMA 350.\\
Fonte: https://www.cjf.jus.br/phpdoc/virtus/
\end{document}
